\lecture{6}{Thu 25 Jan 2024 13:30}{Relative Homology}

\textbf{We want a property in the flavor of}
\[
	\text{``} H_n(X) = H_n(A) \oplus H_n(X - A) \text{''}
.\] 

\begin{definition}[Exact Sequence]
	\label{defn:ExactSequence}
	A \textit{sequence} of abelian groups is a diagram
% https://quiver.local:8080/#q=WzAsNSxbMCwwLCJcXGRvdHMiXSxbMiwwLCJDX3tuKzF9Il0sWzQsMCwiQ19uIl0sWzYsMCwiQ197biAtIDF9Il0sWzgsMCwiXFxkb3RzIl0sWzAsMV0sWzEsMiwiZl97bisxfSJdLFsyLDMsImZfe259Il0sWzMsNF1d
\[\begin{tikzcd}
	\dots && {C_{n+1}} && {C_n} && {C_{n - 1}} && \dots
	\arrow[from=1-1, to=1-3]
	\arrow["{f_{n+1}}", from=1-3, to=1-5]
	\arrow["{f_{n}}", from=1-5, to=1-7]
	\arrow[from=1-7, to=1-9]
\end{tikzcd}\]
such that all compositions are zero.

We remark that a sequence may be bounded; i.e. it may terminate at some index. A chain complex is thus an unbounded sequence.

A sequence $C_*$ is \textit{exact} at degree $n$, if $\operatorname{Im} f_{n + 1} = \operatorname{ker} f_n$.

A sequence $C_*$ is exact if it is exact at every degree $n$.
\end{definition}

\begin{example}
	\begin{itemize}
		\item $0 \to  A \xrightarrow{f} B$ is exact iff $f$ is injective.
		\item $B \xrightarrow{g} C \to  0$ is exact iff $g$ is surjective.
		\item $f$ is injective and  $g$ is surjective iff
% https://quiver.local:8080/#q=WzAsNSxbMCwwLCIwIl0sWzIsMCwiQSJdLFs0LDAsIkIiXSxbNiwwLCJDIl0sWzgsMCwiMCJdLFswLDFdLFsxLDIsImYiXSxbMiwzLCJnIl0sWzMsNF1d
\[\begin{tikzcd}
	0 && A && B && C && 0
	\arrow[from=1-1, to=1-3]
	\arrow["f", from=1-3, to=1-5]
	\arrow["g", from=1-5, to=1-7]
	\arrow[from=1-7, to=1-9]
\end{tikzcd}\]
 is exact. This is called a \textit{short exact sequence.}

 If the sequence is exact, then $B / \operatorname{Im} \cong C$.
	\end{itemize}
\end{example}

\begin{example}
	Let $Z_n = \mathbb{Z} / n \mathbb{Z}$. Consider
% https://quiver.local:8080/#q=WzAsNSxbMCwwLCIwIl0sWzIsMCwiXFxtYXRoYmJ7Wn1fMiJdLFs0LDAsIlxcbWF0aGJie1p9XzQiXSxbNiwwLCJcXG1hdGhiYntafV8yIl0sWzgsMCwiMCJdLFswLDFdLFsxLDIsImYiXSxbMiwzLCJnIl0sWzMsNF1d
\[\begin{tikzcd}
	0 && {\mathbb{Z}_2} && {\mathbb{Z}_4} && {\mathbb{Z}_2} && 0
	\arrow[from=1-1, to=1-3]
	\arrow["f", from=1-3, to=1-5]
	\arrow["g", from=1-5, to=1-7]
	\arrow[from=1-7, to=1-9]
\end{tikzcd}\]
Define $f: [ a + 2 k] \to [2 a + 4 k]$, and $g: [b] \to  [b]$.
\end{example}

Define the category $\text{\textbf{Top}}_2$ such that the objects are pairs  $(X, A)$ where $A$ is a subspace of $X$ ; a morphism $(X, A) \to  (Y, B)$ is a continuous map $f: X \to  Y$ such that $f(A) \subset B$.

\begin{definition}[Relative Chain Group]
	\label{defn:RelativeChagroup}
	Define
	\[
		S_n(X, A) = \frac{S_n(X)}{S_n(A)}
	.\] 
	This is the free abelian group generated by singular $n$-simplices in $X$ that do not lie entirely in $A$.
\end{definition}

Note that: $d(S_n(A)) \subset S_{n - 1}(A)$. So $d$ maps across the quotient:
% https://quiver.local:8080/#q=WzAsNixbMCwwLCJTX24oWCkiXSxbMiwwLCJTX3tuIC0gMX0oWCkiXSxbMCwyLCJTX24oWCwgQSkiXSxbMiwyLCJTX3tuIC0gMX0oWCwgQSkiXSxbMCwzLCJbXFxzaWdtYV0iXSxbMiwzLCJbZChcXHNpZ21hKV0iXSxbMCwxLCJkIl0sWzAsMl0sWzIsMywiZCIsMix7InN0eWxlIjp7ImJvZHkiOnsibmFtZSI6ImRhc2hlZCJ9fX1dLFsxLDNdLFs0LDUsIiIsMix7InNob3J0ZW4iOnsic291cmNlIjozMCwidGFyZ2V0IjozMH0sInN0eWxlIjp7InRhaWwiOnsibmFtZSI6Im1hcHMgdG8ifX19XV0=
\[\begin{tikzcd}
	{S_n(X)} && {S_{n - 1}(X)} \\
	\\
	{S_n(X, A)} && {S_{n - 1}(X, A)} \\
	{[\sigma]} && {[d(\sigma)]}
	\arrow["d", from=1-1, to=1-3]
	\arrow[from=1-1, to=3-1]
	\arrow["d"', dashed, from=3-1, to=3-3]
	\arrow[from=1-3, to=3-3]
	\arrow[shorten <=19pt, shorten >=19pt, maps to, from=4-1, to=4-3]
\end{tikzcd}\]

Now, the \textit{relative chain complex} is defined as
\[
	S_*(X, A) = \left( \ldots \to S_{n + 1}(X, A) \to S_n(X, A) \to \ldots \right) 
.\] 
$H_n(X, A) := H_n (S_*(X, A))$ is the $n$-th \textit{relative homology group}.
Similarly, $Z_n(X, A) = \operatorname{ker} \left( S_n (X, A) \to  S_{n - 1} (X, A) \right) $ is relative $n$-cycles, $B_n(X, A) = \operatorname{Im} \left( S_{n + 1} (X, A) \to  S_n (X, A) \right) $ is the relative $n$ boundary. Thus,
\[
	H_n(X, A) = Z_n(X, A) / B_n(X, A)
.\] 

What does a relative $n$-cycle look like? 
\[
	[c] \in Z_n(X, A) \iff  d([c]) = [d(c)] = 0 \in  S_{n - 1}(X,A) \iff  d(c) \in S_{n - 1}(A)
.\] 

On the other hand,
\begin{multline}
	[c] \in B_n(X, A) \iff \exists   [e] \in S_{n + 1}(X, A), [d(e)] = [c] \in S_n(X, A) \iff \\
	\exists e \in S_{n + 1}(X), \exists b \in S_n(A), d(e) = c + b \in S_n(X)
.\end{multline} 

Consider $[a] \in H_n(X, A) = Z_n(X, A) / B_n(X, A)$, where $Z_n(X, A) \subset S_n(X, A) = S_n(X) / S_n(A)$. Now we have $a \in Z_n(X, A)$ such that $a = [a']$,  $a' \in S_n(X)$. Thus $[a] = \left[ [a'] \right] $. The first bracket on the rhs stands for equivalence class in relative chain group. The second bracket on the rhs stands for element in the relative homology. For notation simplicity we write $[a] = [a']$.


Now, we have
% https://quiver.local:8080/#q=WzAsMjEsWzAsMSwiMCJdLFsxLDEsIlNfbihBKSJdLFsxLDAsIlxcdmRvdHMiXSxbMywwLCJcXHZkb3RzIl0sWzUsMCwiXFx2ZG90cyJdLFszLDEsIlhfbihYKSJdLFs1LDEsIlNfbihYLCBBKSJdLFs2LDEsIjAiXSxbMSwzLCJTX3tuIC0gMX0oQSkiXSxbMywzLCJTX3tuIC0gMX0oWCkiXSxbNSwzLCJTX3tuIC0gMX0oWCwgQSkiXSxbNiwzLCIwIl0sWzAsMywiMCJdLFsxLDQsIjAiXSxbMyw0LCIwIl0sWzUsNCwiMCJdLFszLDJdLFszLDUsIlxcc2lnbWEiXSxbNCw1LCJbXFxzaWdtYV0iXSxbMyw2LCJkIFxcc2lnbWEiXSxbNCw2LCJbIGQgXFxzaWdtYV0iXSxbMCwxXSxbMiwxXSxbMSw1LCJcXGlvdGEiXSxbNSw2LCJcXHBpIl0sWzMsNV0sWzYsN10sWzEsOCwiZCIsMl0sWzksMTAsIlxccGkiXSxbNSw5LCJkIiwyXSxbNiwxMCwiZCJdLFsxMCwxMV0sWzEyLDhdLFs4LDksIlxcaW90YSJdLFs4LDEzXSxbOSwxNF0sWzEwLDE1XSxbNCw2XSxbMTcsNSwiIiwwLHsiY3VydmUiOi00LCJzdHlsZSI6eyJib2R5Ijp7Im5hbWUiOiJkb3R0ZWQifSwiaGVhZCI6eyJuYW1lIjoibm9uZSJ9fX1dLFsxNywxOCwiIiwyLHsic3R5bGUiOnsidGFpbCI6eyJuYW1lIjoibWFwcyB0byJ9fX1dLFsxNywxOSwiIiwyLHsic3R5bGUiOnsidGFpbCI6eyJuYW1lIjoibWFwcyB0byJ9fX1dLFsxOSwyMCwiIiwyLHsic3R5bGUiOnsidGFpbCI6eyJuYW1lIjoibWFwcyB0byJ9fX1dLFsxOCwyMCwiIiwxLHsic3R5bGUiOnsidGFpbCI6eyJuYW1lIjoibWFwcyB0byJ9fX1dXQ==
\[\begin{tikzcd}
	& \vdots && \vdots && \vdots \\
	0 & {S_n(A)} && {X_n(X)} && {S_n(X, A)} & 0 \\
	&&& {} \\
	0 & {S_{n - 1}(A)} && {S_{n - 1}(X)} && {S_{n - 1}(X, A)} & 0 \\
	& 0 && 0 && 0 \\
	&&& \sigma & {[\sigma]} \\
	&&& {d \sigma} & {[ d \sigma]}
	\arrow[from=2-1, to=2-2]
	\arrow[from=1-2, to=2-2]
	\arrow["\iota", from=2-2, to=2-4]
	\arrow["\pi", from=2-4, to=2-6]
	\arrow[from=1-4, to=2-4]
	\arrow[from=2-6, to=2-7]
	\arrow["d"', from=2-2, to=4-2]
	\arrow["\pi", from=4-4, to=4-6]
	\arrow["d"', from=2-4, to=4-4]
	\arrow["d", from=2-6, to=4-6]
	\arrow[from=4-6, to=4-7]
	\arrow[from=4-1, to=4-2]
	\arrow["\iota", from=4-2, to=4-4]
	\arrow[from=4-2, to=5-2]
	\arrow[from=4-4, to=5-4]
	\arrow[from=4-6, to=5-6]
	\arrow[from=1-6, to=2-6]
	\arrow[curve={height=-24pt}, dotted, no head, from=6-4, to=2-4]
	\arrow[maps to, from=6-4, to=6-5]
	\arrow[maps to, from=6-4, to=7-4]
	\arrow[maps to, from=7-4, to=7-5]
	\arrow[maps to, from=6-5, to=7-5]
\end{tikzcd}\]

\begin{definition}[Short Exact Sequence of Chain Complexes]
	\label{defn:ShortExactSequenceChainComplexes}
	
	A s.e.s. of chain complexes is $0 \to A_* \xrightarrow{f} B_* \xrightarrow{g} C_* \to  0$, where $f, g$ are chain maps, and for each $n$, $0 \to A_n \xrightarrow{f} B_n \xrightarrow{g} C_n \to  0$ is a s.e.s.
\end{definition}

This induces an exact sequence of homologies (not in general short exact):

\begin{lemma}
	\label{lem:ShortShortSequenceHomology}
	The short exact sequence of chain complexes  $0 \to A_* \xrightarrow{f} B_* \xrightarrow{g} C_* \to  0$ induces an exact sequence on homologies:
% https://quiver.local:8080/#q=WzAsNyxbMCwwLCJIX24oQV8qKSJdLFsyLDAsIkhfbihCXyopIl0sWzQsMCwiSF9uKENfKikiXSxbMCwxLCJbYV9uXSJdLFsyLDEsIltmX24gKGFfbildIl0sWzIsMiwiW2Jfbl0iXSxbNCwyLCJbZ19uKGJfbildIl0sWzAsMSwiZl8qIl0sWzEsMiwiZ18qIl0sWzMsNCwiIiwwLHsic2hvcnRlbiI6eyJzb3VyY2UiOjIwLCJ0YXJnZXQiOjIwfSwic3R5bGUiOnsidGFpbCI6eyJuYW1lIjoibWFwcyB0byJ9fX1dLFs1LDYsIiIsMCx7InNob3J0ZW4iOnsic291cmNlIjoyMCwidGFyZ2V0IjoyMH0sInN0eWxlIjp7InRhaWwiOnsibmFtZSI6Im1hcHMgdG8ifX19XV0=
\[\begin{tikzcd}
	{H_n(A_*)} && {H_n(B_*)} && {H_n(C_*)} \\
	{[a_n]} && {[f_n (a_n)]} \\
	&& {[b_n]} && {[g_n(b_n)]}
	\arrow["{f_*}", from=1-1, to=1-3]
	\arrow["{g_*}", from=1-3, to=1-5]
	\arrow[shorten <=10pt, shorten >=10pt, maps to, from=2-1, to=2-3]
	\arrow[shorten <=10pt, shorten >=10pt, maps to, from=3-3, to=3-5]
\end{tikzcd}\]
\end{lemma}
\begin{proof}
Since $gf = 0$, we have $g_* f_* = 0$, so $\operatorname{Im} f_* \subset \operatorname{ker} g_*$.

We check for $\operatorname{ker} g_* \subset \operatorname{Im} f_*$. Take $[b_n] \in \operatorname{ker} g_*$, $b_n \in Z_n(B_*)$. Then $g_*([b_n]) = [g(b_n)] = c_n = 0 \in H_n(C_*)$, which means $g(b_n) \in B_n(C_*)$. Thus, there exists $c_{n + 1} \in C_{n + 1}$, $d (c_{n + 1}) = c_n = g(b_n)$. Since $g$ is surjective, there exists $b_{n + 1} \in B_{n + 1}$, $g(b_{n + 1}) = c_{n + 1}$. So now we have $ d\circ  g_{n + 1} (b_{n + 1} ) = g_{n} (b_n) = c_n \in B_n(C)$.

Set $b_n' = d b_{n + 1}$, since diagram commutes, we have $d g_{n + 1} = g_n d$. Thus $g_n(b_n') = c_n$, so $g_n(b_n - b_n') = 0$. This means $(b_n - b_n') \in \operatorname{ker} g_n = \operatorname{Im} f_n$. But the horizontal chain is exact, i.e. $\operatorname{ker} g_n = \operatorname{Im} f_n$, so there exists $a_n \in A_n$, $f_n(a_n) = b_n - b_n' = b_n - d(b_{n + 1}).$ We want to show $a_n$ is a cycle. But $f_{n - 1} d(a_n) = d(f_n(a_n)) = 0$, so  $d(a_n) = 0$, since $f_{n - 1}$ is injective. 

Thus, $f_*[a_n] = [b_n - b_n'] = [b_n] \in H_n(B_*)$.

Use the following diagrams to guide thinking:
% https://quiver.local:8080/#q=WzAsMTMsWzIsMCwiXFxkb3RzIl0sWzQsMCwiQl97bisxfSJdLFs2LDAsIkNfe24gKyAxfSJdLFs4LDAsIjAiXSxbNCwyLCJCX24iXSxbNiwyLCJDX24iXSxbMiwyLCJBX24iXSxbMCwyLCIwIl0sWzgsMiwiMCJdLFsyLDQsIkFfe24gLSAxfSJdLFs0LDQsIkJfe24gLSAxfSJdLFswLDQsIjAiXSxbNiw0LCJcXGRvdHMiXSxbMCwxXSxbMSwyLCJnIiwyXSxbMiwzXSxbNyw2XSxbNiw0LCJmIiwyXSxbNCw1LCJnIiwyXSxbNSw4XSxbMTEsOV0sWzksMTAsImYiLDJdLFsxMCwxMl0sWzEsNCwiZCIsMV0sWzIsNSwiZCIsMV0sWzYsOSwiZCIsMV0sWzQsMTAsImQiLDFdLFs0LDUsIjEiLDAseyJvZmZzZXQiOi0yLCJjb2xvdXIiOlsxMjAsNjAsNjBdLCJzdHlsZSI6eyJib2R5Ijp7Im5hbWUiOiJzcXVpZ2dseSJ9fX0sWzEyMCw2MCw2MCwxXV0sWzUsMiwiMiIsMCx7Im9mZnNldCI6LTIsImNvbG91ciI6WzEyMCw2MCw2MF0sInN0eWxlIjp7ImJvZHkiOnsibmFtZSI6InNxdWlnZ2x5In19fSxbMTIwLDYwLDYwLDFdXSxbMiwxLCIzIiwyLHsib2Zmc2V0IjoyLCJjb2xvdXIiOlsxMjAsNjAsNjBdLCJzdHlsZSI6eyJib2R5Ijp7Im5hbWUiOiJzcXVpZ2dseSJ9fX0sWzEyMCw2MCw2MCwxXV0sWzQsNiwiNCIsMix7Im9mZnNldCI6MiwiY29sb3VyIjpbMTIwLDYwLDYwXSwic3R5bGUiOnsiYm9keSI6eyJuYW1lIjoic3F1aWdnbHkifX19LFsxMjAsNjAsNjAsMV1dLFs2LDksIjUiLDIseyJvZmZzZXQiOjMsImNvbG91ciI6WzEyMCw2MCw2MF0sInN0eWxlIjp7ImJvZHkiOnsibmFtZSI6InNxdWlnZ2x5In19fSxbMTIwLDYwLDYwLDFdXV0=
\[\begin{tikzcd}
	&& \dots && {B_{n+1}} && {C_{n + 1}} && 0 \\
	\\
	0 && {A_n} && {B_n} && {C_n} && 0 \\
	\\
	0 && {A_{n - 1}} && {B_{n - 1}} && \dots
	\arrow[from=1-3, to=1-5]
	\arrow["g"', from=1-5, to=1-7]
	\arrow[from=1-7, to=1-9]
	\arrow[from=3-1, to=3-3]
	\arrow["f"', from=3-3, to=3-5]
	\arrow["g"', from=3-5, to=3-7]
	\arrow[from=3-7, to=3-9]
	\arrow[from=5-1, to=5-3]
	\arrow["f"', from=5-3, to=5-5]
	\arrow[from=5-5, to=5-7]
	\arrow["d"{description}, from=1-5, to=3-5]
	\arrow["d"{description}, from=1-7, to=3-7]
	\arrow["d"{description}, from=3-3, to=5-3]
	\arrow["d"{description}, from=3-5, to=5-5]
	\arrow["1", shift left=2, color={rgb,255:red,92;green,214;blue,92}, squiggly, from=3-5, to=3-7]
	\arrow["2", shift left=2, color={rgb,255:red,92;green,214;blue,92}, squiggly, from=3-7, to=1-7]
	\arrow["3"', shift right=2, color={rgb,255:red,92;green,214;blue,92}, squiggly, from=1-7, to=1-5]
	\arrow["4"', shift right=2, color={rgb,255:red,92;green,214;blue,92}, squiggly, from=3-5, to=3-3]
	\arrow["5"', shift right=3, color={rgb,255:red,92;green,214;blue,92}, squiggly, from=3-3, to=5-3]
\end{tikzcd}\]

Remark: the proof proceeds in two parts. For the first part, we want to transform the statement $[b_n] \in \operatorname{ker} g_*$ to something about $b_n$, so that we can use the exactness of $0\to A_n \to B_n \to C_n \to 0$. This is done by showing that there exists some $b_n' \in B_n(B_*)$ such that $b_n - b_n' \in \operatorname{ker} g$, and $d(b_n - b_n') = 0$.

The second part is using the exactness property to obtain some $a_n$ such that $f(a_n) = b_n - b_n'$, and $a_n$ happens to be a cycle, so we can verify $f_*([a_n]) = [b_n - b_n'] = [b_n]$.
\end{proof}

This is not in general a short exact sequence; it is exact in the middle, but not in general exact on the sides. However, we can connect them into a long exact sequence, by connecting $H_n(C_*) \to  H_{n - 1}(A_*)$.
